\subsection{Синхронизация и синфазность}

Общим наблюдением среди всех видов связи является присутствие областей частичной
синхронизации осцилляторов: языков захвата пары $0\leftrightarrow 1$ и пары
$0\leftrightarrow 2$, а также тривиальной линии с равными расстройками
$\delta_1 = \delta_2$ (рисунок~\ref{fig:regimes-grid}). С увеличением силы связи языки
становятся шире и захватывают всё большую часть пространства параметров, а на их
пересечении возникает область полной синхронизации.

Режим на всех картах определяется как мода по 30 начальным условиям. Поэтому карту
режимов всегда следует читать вместе с картой мультистабильности
(рисунок~\ref{fig:multi-grid}): там, где доля начальных условий вне доминирующего режима
велика, мода скрывает расхождение запусков. Отметим, что мода "--- это самое частое
значение, а не большинство: при 30 запусках и трёх равных по частоте кодах доля вне моды
достигает $2/3$, что и определяет верхнюю границу шкалы.

Карты синфазности (рисунок~\ref{fig:inphase-grid}) показывают, что частотный захват и
синфазность "--- не одно и то же, причём у двух типов связи они ведут себя противоположно.
При диссипативной связи область полной синхронизации отвечает повышенной синфазности:
при $\varepsilon = 0.05$ медиана $P$ составляет там $0.763$ против $0.500$ в незахваченной
части карты. При инерционной связи, напротив, полная синхронизация сопровождается
\emph{пониженной} синфазностью: медиана $P$ равна $0.445$ без нормировки и $0.369$ с
нормировкой, то есть ниже, чем в отсутствие захвата ($0.498$).

Из этого следует, что инерционная связь захватывает частоты, но разводит фазы, тогда как
диссипативная сводит их вместе. Синхронизация оказывается ещё и синфазной только при
диссипативной связи.

Линия $\delta_1 = \delta_2$ отдельного языка не образует. На ней листья идентичны, и
захват пары $1\leftrightarrow 2$ следует из перестановочной симметрии сети, а не из
взаимодействия листьев между собой; вдали от пересечения языков её ширина по расстройке
не превышает шага сетки ни при какой силе связи. Вблизи центра
(рисунок~\ref{fig:regimes-grid-zoom}) полоса $1\leftrightarrow 2$ заметно шире, однако это не
признак самостоятельного языка: там оба листа близки к захвату с хабом, их наблюдаемые
частоты $\Omega_{01}$ и $\Omega_{02}$ малы, поэтому разность
$\Omega_{12} = \Omega_{02} - \Omega_{01}$ оказывается меньше порога уже на целой двумерной
области.

\begin{figure}[htbp]
  \centering
  \includegraphics[width=\textwidth]{plots/regimes-grid.png}
  \caption{Распределение режимов частотной синхронизации на плоскости расстроек
           $(\delta_1, \delta_2)$ для четырёх вариантов связи. Цветом показан режим,
           наиболее часто реализующийся среди 30 начальных условий при данном наборе
           параметров}
  \label{fig:regimes-grid}
\end{figure}

\begin{figure}[htbp]
  \centering
  \includegraphics[width=\textwidth]{plots/regimes-grid-zoom.png}
  \caption{Центральная область карт режимов частотной синхронизации на плоскости
           расстроек $(\delta_1, \delta_2)$. Увеличенный фрагмент показывает положение
           языков захвата в окрестности линии $\delta_1 = \delta_2$}
  \label{fig:regimes-grid-zoom}
\end{figure}

\begin{figure}[htbp]
  \centering
  \includegraphics[width=\textwidth]{plots/multi-grid.png}
  \caption{Оценка мультистабильности на плоскости расстроек $(\delta_1, \delta_2)$.
           Цветом показана доля начальных условий, для которых реализуется режим,
           отличный от наиболее часто наблюдаемого режима при тех же параметрах}
  \label{fig:multi-grid}
\end{figure}

\begin{figure}[htbp]
  \centering
  \includegraphics[width=\textwidth]{plots/inphase-grid.png}
  \caption{Синфазность осцилляторов на плоскости расстроек $(\delta_1, \delta_2)$.
           Цветом показана медиана показателя синфазности $P$, вычисленная по 30
           начальным условиям}
  \label{fig:inphase-grid}
\end{figure}

\subsubsection{Инерционная связь}

На рисунке~\ref{fig:regimes-grid-zoom} видно, что при ненормированной инерционной
связи языки захвата центрированы не в нуле, а на $\delta \approx \varepsilon/2$. Причина
в том, что ненормированная связь через оператор Лапласа содержит слагаемое
$-\varepsilon k_n x_n$, то есть добавляет к возвращающей силе вклад, пропорциональный
степени узла. Эффективная частота узла оказывается равной
$\omega^2_{\text{eff}} = \omega_n^2 + \varepsilon k_n$. У хаба $k_0 = 2$, у листьев
$k_1 = k_2 = 1$, поэтому хаб смещается по частоте вдвое сильнее, и для захвата лист должен
быть заранее расстроен вверх примерно на $\varepsilon/2$. Нормировка на степень узла
выравнивает эту добавку, и языки оказываются центрированы на нуле.

Различие между двумя вариантами инерционной связи наиболее наглядно проявляется при
$\varepsilon = 0.1$ (рисунок~\ref{fig:special-inertial}). При ненормированной связи полная
синхронизация занимает $12.1\%$ ячеек центральной области, однако $63.2\%$ ячеек этой
области мультистабильны, а медианная доля начальных условий вне доминирующего режима
составляет $0.167$. Из этого следует, что полная синхронизация в центре устойчиво не
устанавливается: при одних и тех же параметрах часть начальных условий приводит к ней, а
часть "--- нет. При нормированной связи полная синхронизация занимает $45.5\%$ центральной
области, а медианная доля вне доминирующего режима равна нулю, то есть режим
устанавливается независимо от выбора начальных условий. Разрушения полной синхронизации
здесь не происходит.

\begin{figure}[htbp]
  \centering
  \includegraphics[width=\textwidth]{plots/special-inertial.png}
  \caption{Влияние нормировки инерционной связи на структуру режимов в центральной
           области плоскости расстроек при $\varepsilon = 0.1$. Сопоставлены карты
           режима синхронизации и доли начальных условий вне доминирующего режима для
           ненормированной и нормированной связи}
  \label{fig:special-inertial}
\end{figure}

\subsubsection{Диссипативная связь}

Диссипативная связь даёт заметно более широкие языки захвата при тех же значениях
$\varepsilon$: уже при $\varepsilon = 0.05$ область без синхронизации занимает $53.7\%$
карты против $90.8\%$ у инерционной связи (рисунок~\ref{fig:regimes-grid}).

Дальнейшее увеличение силы связи выявляет существенное различие между нормированным и
ненормированным вариантами (рисунок~\ref{fig:special-dissipative}). При $\varepsilon = 0.07$
у ненормированной связи область без синхронизации занимает всего $5.2\%$ карты, то есть
захват наблюдается практически во всём исследованном диапазоне расстроек. Нормированная
связь при том же $\varepsilon$ такой широкой зоны захвата не даёт: область без
синхронизации занимает $75.9\%$ карты.

При переходе к $\varepsilon = 0.1$ у нормированной связи наблюдается слияние языков
частичного захвата в область полной синхронизации. Доля полной синхронизации возрастает с
$2.5\%$ до $13.8\%$, тогда как доли частичных режимов $0\leftrightarrow 1$ и
$0\leftrightarrow 2$ при этом слегка уменьшаются, с $10.6\%$ до $9.8\%$. Из этого следует,
что область полной синхронизации разрастается не в незахваченную часть пространства
параметров, а поглощает окрестность пересечения языков.

\begin{figure}[htbp]
  \centering
  \includegraphics[width=\textwidth]{plots/special-dissipative.png}
  \caption{Изменение областей частотного захвата при диссипативной связи. Карты
           показывают расширение языков частичной синхронизации и формирование области
           полной синхронизации при увеличении силы связи}
  \label{fig:special-dissipative}
\end{figure}
